% !TEX root = ../main.tex
\chapter{不等式约束与 KKT 条件}
\label{ch:kkt}

\noindent\emph{用到的前置工具：梯度与 Jacobian（多元微分一章，第~\ref{ch:multideriv}~章）；等式约束下的 Lagrange 乘子法（见前一章「等式约束与 Lagrange 乘子」）；凹凸函数与凸集（第~\ref{ch:convex}~章）。}
\medskip

上一章的 Lagrange 乘子法处理的是``恰好取等''的约束：预算\emph{花光}、资源\emph{用尽}、产量\emph{正好}等于订单。可现实里的约束多半不是这副样子。消费者``花费不超过收入''——但他完全可以存一点；厂商``产量不超过产能''——但淡季时产能根本闲着；几乎每个经济变量还自带一条``不能为负''的底线。这些都是\textbf{不等式}约束 $g(\vc x)\le 0$，它带来一个 Lagrange 法没有的新问题：约束到底起不起作用，要等解出来才知道。本章要讲的 KKT 条件，正是把 Lagrange 法推广到不等式约束的那把工具——它用一组优雅的``互补''关系，一次性地把``哪些约束紧、哪些松''和``各自的乘子''都刻画了出来。这是带约束优化真正的主力一阶条件，也是凸优化、受约束估计、机制设计里反复登场的同一套语言。

\section{不等式约束带来的新问题}
\label{sec:kkt-1}

先看等式约束法在哪里不够用。考虑标准的约束极大化问题
\[
    \max_{\vc x\in\RR^{n}}\; f(\vc x)
    \qquad\text{s.t.}\qquad
    g_j(\vc x)\le 0,\quad j=1,\dots,m .
\]
（任何``$\ge$''约束写成 $-g\le0$，等式约束 $h=0$ 拆成 $h\le0$ 与 $-h\le0$ 即可，所以这一种形式已经足够一般。）

如果把每条 $g_j\le0$ 都当成 $g_j=0$ 去套 Lagrange 法，立刻会出错。问题在于：在最优点处，有些约束\emph{恰好取等}（$g_j(\vc x^*)=0$，我们说这条约束是\textbf{起作用的}或\textbf{紧的}，active / binding），有些则\emph{严格松弛}（$g_j(\vc x^*)<0$，约束是\textbf{不起作用的}或\textbf{松的}，inactive / slack）。

\state{不等式约束的核心难点：起作用集是内生的}{
紧的约束像一堵墙，最优点被它顶住，它会像等式约束那样在一阶条件里贡献一个乘子；松的约束则形同虚设，最优点离它还远，挪一挪它丝毫不影响最优，它\emph{不该}出现在一阶条件里。麻烦在于：\textbf{究竟哪些约束紧、哪些松，是问题的解的一部分，事先并不知道}。我们不能像等式约束那样直接列方程，必须让一阶条件本身有能力``自动判别''每条约束的死活——这正是 KKT 条件比 Lagrange 法多出来的那层结构。
}

\begin{figure}[h]
\centering
\begin{tikzpicture}[scale=1.0,>=Stealth]
    % 坐标轴
    \draw[->] (-0.2,0) -- (6.2,0) node[right] {$x_1$};
    \draw[->] (0,-0.2) -- (0,6.2) node[above] {$x_2$};
    % 可行域：x1>=0, x2>=0, x1+x2<=5。三角形 (0,0)(5,0)(0,5)
    \fill[blue!8] (0,0) -- (5,0) -- (0,5) -- cycle;
    \draw[thick,blue!60!black] (5,0) -- (0,5);
    % 约束标签（可行域内右下空白处，远离圆与梯度）
    \node[blue!60!black,anchor=west] at (2.55,1.15) {$g_1=x_1+x_2-5\le 0$};
    % 无约束最优点 (4,4)，在可行域外
    \fill[black] (4,4) circle (1.6pt);
    \node[anchor=west] at (4.12,3.92) {$\vc x^{\mathrm{u}}=(4,4)$};
    \node[anchor=west] at (4.12,4.45) {无约束最优};
    % 目标等值线：以(4,4)为圆心的圆（f=-((x1-4)^2+(x2-4)^2)）
    % 半径选 r=2.121=3/sqrt2 时恰与可行域相切于(2.5,2.5)
    \draw[gray!70,thin] (4,4) circle (1.0);
    \draw[gray!70,thin] (4,4) circle (1.6);
    \draw[red!70!black,thick] (4,4) circle (2.1213); % 相切圆
    % 约束极大点 (2.5,2.5)
    \fill[black] (2.5,2.5) circle (1.8pt);
    \node[anchor=north east] at (2.42,2.40) {$\vc x^{*}$};
    % 等值线方向标注（左下空白）
    \node[red!70!black,anchor=west] at (0.10,0.55) {目标等值线 $f=$ 常数};
    \draw[red!70!black,->] (1.05,0.85) -- (1.95,1.85); % 引线指向相切圆左下弧，不穿其它曲线
    % 两个梯度箭头：grad f 与 grad g1 共线（外法向 (1,1)方向）。
    % 为可读起见把二者沿垂直方向(-1,1)略微错开，仍保持平行以示共线。
    % grad f 偏下侧 (0.13,-0.13)，grad g1 偏上侧 (-0.13,0.13)
    \draw[->,thick,red!70!black] (2.63,2.37) -- (3.13,2.87);
    \node[red!70!black,anchor=north west] at (3.15,2.80) {$\grad f$};
    \draw[->,thick,blue!60!black] (2.37,2.63) -- (3.37,3.63);
    \node[blue!60!black,anchor=south east] at (3.34,3.66) {$\grad g_1$};
\end{tikzpicture}
\caption{不等式约束极大化的几何。可行域（浅蓝三角）由 $x_1\ge0,\,x_2\ge0,\,x_1+x_2\le5$ 围成；目标 $f$ 的等值线是以无约束最优 $\vc x^{\mathrm u}=(4,4)$ 为圆心、半径越小值越大的同心圆。无约束最优落在可行域外，于是约束最优 $\vc x^*=(2.5,2.5)$ 被边界 $g_1=0$ 顶住：等值圆在此与边界相切，$\grad f$ 与起作用约束的 $\grad g_1$ 同向（都指向可行域外侧），二者共线——这正是 KKT 驻点条件 $\grad f=\lambda_1\grad g_1$（$\lambda_1>0$）的图像。$x_1\ge0,\,x_2\ge0$ 两条约束在此处松弛，对应乘子为零。}
\label{fig:kkt-1}
\end{figure}

图~\ref{fig:kkt-1} 把这件事画了出来。当无约束最优落在可行域之外，最优解被某条边界顶住，在那里目标的等值线与边界相切，目标梯度 $\grad f$ 沿着``往可行域外推''的方向，恰好与起作用约束的梯度 $\grad g_j$ 共线。下面要做的，是把这幅图翻译成一组代数条件。

\section{KKT 条件：四件套}
\label{sec:kkt-2}

仿照 Lagrange 法，给每条约束 $g_j\le0$ 配一个乘子 $\lambda_j$，写出 Lagrange 函数
\[
    \cL(\vc x,\vc\lambda)
    = f(\vc x)-\sum_{j=1}^{m}\lambda_j\,g_j(\vc x),
    \qquad \vc\lambda=(\lambda_1,\dots,\lambda_m) .
\]
（这里取减号，是为了让 $\lambda_j\ge0$；符号约定的来由稍后细说。）刻画最优点的，是下面四组条件。

\thm{Karush--Kuhn--Tucker 必要条件}{
设 $f,g_1,\dots,g_m$ 在 $\vc x^*$ 附近连续可微，$\vc x^*$ 是上述问题的局部极大点，且在 $\vc x^*$ 处满足某个\emph{约束规范}（见第~\ref{sec:kkt-cq}~节）。则存在乘子 $\lambda_1^*,\dots,\lambda_m^*$，使得：
\begin{align*}
    &\text{(i) 驻点（stationarity）：}
    && \grad f(\vc x^*)=\sum_{j=1}^{m}\lambda_j^*\,\grad g_j(\vc x^*); \\[2pt]
    &\text{(ii) 原始可行（primal feasibility）：}
    && g_j(\vc x^*)\le 0,\quad j=1,\dots,m; \\[2pt]
    &\text{(iii) 对偶可行（dual feasibility）：}
    && \lambda_j^*\ge 0,\quad j=1,\dots,m; \\[2pt]
    &\text{(iv) 互补松弛（complementary slackness）：}
    && \lambda_j^*\,g_j(\vc x^*)=0,\quad j=1,\dots,m .
\end{align*}
}

这四条各管一摊，缺一不可，值得逐条体会其含义。

\paragraph{(i) 驻点。}
这是 Lagrange 一阶条件的直接翻版：在最优点，目标梯度被起作用约束的梯度``撑住''。写开就是
\[
    \grad_{\vc x}\cL(\vc x^*,\vc\lambda^*)
    =\grad f(\vc x^*)-\sum_{j}\lambda_j^*\,\grad g_j(\vc x^*)=\vc 0 .
\]
它说的是：在最优点再没有一个``既可行又能让 $f$ 上升''的移动方向了。任何想让 $f$ 增大的方向（即与 $\grad f$ 夹角小于 $90^\circ$ 的方向），都会顶到某条起作用约束的墙上。

\paragraph{(ii) 原始可行。}
平凡但不可省：最优点首先得在可行域里。

\paragraph{(iii) 对偶可行。}
乘子非负，$\lambda_j^*\ge0$。这是不等式约束相对等式约束最关键的新增条件——等式约束的乘子可正可负，不等式约束的乘子却被钉死在非负一侧。其根源是不等式约束\emph{只朝一个方向}限制你（见下节符号约定）。

\paragraph{(iv) 互补松弛。}
$\lambda_j^*\,g_j(\vc x^*)=0$ 对每个 $j$ 成立——这是 KKT 的灵魂。它说：对每条约束，乘子 $\lambda_j^*$ 与松弛量 $g_j(\vc x^*)$ 至少有一个为零，二者不能同时非零。换言之

\state{互补松弛把``起作用集是内生的''这一难点一举化解}{
互补松弛 $\lambda_j^* g_j(\vc x^*)=0$ 等价于一句大白话：
\begin{itemize}
    \item 约束\emph{松}（$g_j(\vc x^*)<0$，离墙还远）$\To$ 它的乘子 $\lambda_j^*=0$，它不进一阶条件，仿佛不存在；
    \item 乘子\emph{为正}（$\lambda_j^*>0$，这条约束在出力）$\To$ 必有 $g_j(\vc x^*)=0$，约束\emph{紧}，最优点正贴在这堵墙上。
\end{itemize}
于是我们不必事先知道哪些约束起作用：互补松弛让一阶条件\emph{自动}关掉所有松约束、只留下紧约束，恰好实现了第~\ref{sec:kkt-1}~节那个``自动判别约束死活''的诉求。乘子 $\lambda_j^*$ 在这里有了清晰的经济含义——它是把第 $j$ 条约束\emph{放松一单位}（资源多给一点、预算多一块钱）所能换来的目标增量，即约束的\textbf{影子价格}。松约束的影子价格自然为零：你手上的资源还没用完，再多给也不值钱。
}

\rmkb[互补松弛不是说``恰好一个为零'']{
$\lambda_j^* g_j=0$ 允许两者\emph{同时}为零：约束恰好取等（$g_j=0$）但其乘子也恰是 $0$。这种``临界''情形（约束刚刚好顶上、却不出力）在分析里要小心，它对应起作用集发生切换的边界，也是后面比较静态可能不光滑的地方。严格互补（每条起作用约束的乘子都 $>0$）是一个更强、更``好''的假设，常被额外加上以保证解的良态。
}

\subsection{符号约定：为什么 $\lambda\ge0$}
\label{sec:kkt-sign}

乘子非负不是约定俗成，而是几何强加的。在一个起作用约束 $g_j(\vc x^*)=0$ 处，可行域局部位于 $g_j\le0$ 一侧，梯度 $\grad g_j$ 指向 $g_j$ 增大的方向——也就是\emph{离开}可行域、朝禁区去的外法向。最优点既然被这堵墙顶住，目标必定还想往墙外走才对：$\grad f$ 必须指向可行域\emph{外侧}，即与外法向 $\grad g_j$ 同向（在多约束时，是落在诸 $\grad g_j$ 张成的``外推锥''里）。``同向''就意味着展开式
\[
    \grad f(\vc x^*)=\sum_{j:\,\text{起作用}}\lambda_j^*\,\grad g_j(\vc x^*)
\]
里的系数 $\lambda_j^*$ 必须非负。若某个 $\lambda_j^*<0$，则 $\grad f$ 有一个指向可行域\emph{内部}的分量，意味着``往里挪一点 $f$ 还能涨''——那 $\vc x^*$ 根本不是极大点，矛盾。

\state{$\lambda\ge0$ 的方向逻辑（一句话记住）}{
不等式约束\textbf{只朝一个方向挡你}：它不让你越过墙，却随你后退。所以在最优点，目标的``上山方向''只可能朝墙外、不可能朝墙里——$\grad f$ 与外法向 $\grad g$ 同侧，乘子因此非负。等式约束两侧都挡，乘子才可正可负。\emph{记号警示}：$\lambda\ge0$ 的正负号死死绑定在你怎么写约束（$g\le0$ 还是 $g\ge0$）和 $\cL$ 里用加号还是减号上；换了写法符号就翻。本书统一取``$g\le0$、$\cL=f-\sum\lambda_j g_j$、$\lambda\ge0$''这一套，用前务必核对自己的约定是否一致。
}

\subsection{从 Lagrange 到 KKT 的推导直觉}
\label{sec:kkt-derive}

KKT 条件可以这样``凑''出来：在最优点 $\vc x^*$，把起作用集记为 $A=\set{j:g_j(\vc x^*)=0}$。只有起作用的约束才在局部约束你，松约束因为有余量，在 $\vc x^*$ 的一个小邻域内自动满足、可以暂时忽略。于是\emph{在该邻域内}，问题退化成一个只带等式约束 $\set{g_j=0:j\in A}$ 的 Lagrange 问题，其一阶条件正是
\[
    \grad f(\vc x^*)=\sum_{j\in A}\lambda_j^*\,\grad g_j(\vc x^*) .
\]
对松约束 $j\notin A$，令 $\lambda_j^*=0$，上式就能写成对所有 $j$ 求和的形式 (i)；而``$j\notin A\To\lambda_j^*=0$''与``$j\in A\To g_j=0$''合起来恰是互补松弛 (iv)。剩下的非负性 (iii) 由上一小节的方向论证补足，可行性 (ii) 是最优点的定义。四件套就此凑齐。

这套``先猜起作用集、再当等式约束解''的思路，也正是手算 KKT 问题的标准套路（见下例）：枚举若干种起作用集的情形，每种情形下把紧约束当等式解出候选点，再回头验证它满足全部四条 KKT 条件——尤其是被忽略约束的可行性与所求乘子的非负性。

\rmkb[严格的推导需要约束规范]{
上面``退化成等式约束 Lagrange 问题''这一步偷偷用了一个前提：起作用约束的梯度 $\set{\grad g_j(\vc x^*):j\in A}$ 表现得足够``规整''，使得可行域在 $\vc x^*$ 处的局部形状真的被这些梯度的线性结构刻画。这个前提就是\emph{约束规范}。它不成立时，乘子可能根本不存在，KKT 条件随之失效——下节专门讨论。严格的证明（用 Farkas 引理或分离超平面把``不存在可行上升方向''翻译成``$\grad f$ 落在 $\grad g_j$ 张成的锥里''）属于较重的凸分析构造，这里只给上述直觉，不逐行搭建。
}

\section{约束规范：为什么需要它}
\label{sec:kkt-cq}

KKT 是\emph{必要}条件，但有一个前提常被初学者忽略：它只在最优点满足某个\textbf{约束规范}（constraint qualification, CQ）时才保证乘子存在。约束规范是对可行域在最优点处``不太病态''的技术性要求；它失效时，一个货真价实的最优点也可能凑不出满足 (i)--(iv) 的乘子。

\defn{线性无关约束规范（LICQ）}{
称问题在可行点 $\vc x^*$ 处满足\textbf{线性无关约束规范}，如果起作用约束的梯度族
\[
    \set{\grad g_j(\vc x^*):j\in A},\qquad A=\set{j:g_j(\vc x^*)=0}
\]
线性无关。
}

LICQ 是最常用、也最容易检验的一种约束规范。此外还有更弱的若干种，例如 Mangasarian \mbox{--}\,Fromovitz 条件与 Slater 条件；凸问题下尤以后者常用（见下节）。它的作用是保证：可行域在 $\vc x^*$ 处的``可行方向''真的由这些梯度的线性结构决定，从而上一节那套``退化成等式 Lagrange''的推导站得住脚。

CQ 为何不可省？根子在于 KKT 驻点条件 (i) 谈的是\emph{梯度}之间的线性关系，而梯度只看约束函数的一阶信息。如果可行域的真实形状在 $\vc x^*$ 处出现了一阶信息``看不见''的尖角或退化（比如两条约束在此相切、梯度共线），那么``$\grad f$ 落在 $\grad g_j$ 张成的锥里''这件事就可能对真正的最优点也不成立。

\ex[约束规范失效，KKT 凑不出乘子]{
求解 $\displaystyle\max_{x_1,x_2}\;x_1$，约束为
\[
    g_1(\vc x)=x_2-(1-x_1)^3\le 0,\qquad g_2(\vc x)=-x_2\le 0 .
\]
可行域被 $x_2\ge0$ 与 $x_2\le(1-x_1)^3$ 夹住。直观上 $x_1$ 越大越好，最优点是 $\vc x^*=(1,0)$（此处两约束都取等）。检验 KKT 是否成立。
}
\sol{
在 $\vc x^*=(1,0)$ 处两约束都起作用。算梯度：
\[
    \grad g_1=\pr{3(1-x_1)^2,\;1}\Big|_{(1,0)}=(0,1),
    \qquad
    \grad g_2=(0,-1) .
\]
两者共线（一个是另一个的 $-1$ 倍），\textbf{线性相关}——LICQ 在此失效。再看 KKT 驻点条件需要
\[
    \grad f=(1,0)=\lambda_1(0,1)+\lambda_2(0,-1)=(0,\lambda_1-\lambda_2) .
\]
左边第一个分量是 $1$，右边第一个分量恒为 $0$：方程无解。于是\textbf{在真正的最优点 $(1,0)$ 处根本不存在 KKT 乘子}。这不是因为 $(1,0)$ 不是最优，而是因为约束规范失效——可行域在该点收成一个``尖''（两条边界相切），一阶梯度信息看不出 $x_1$ 还能不能增大。这个反例提醒我们：套用 KKT 前，至少要瞥一眼起作用约束的梯度是不是线性无关。
}

\state{约束规范的实用态度}{
绝大多数应用里，约束要么是\emph{线性}的（预算、资源、非负约束——线性约束自动满足一种约束规范，无需另检），要么可行域``凸而内点非空''（满足 Slater 条件，见下节）。这两类情形覆盖了经济学里的几乎全部问题，所以约束规范通常不是障碍。真正要警惕的，是非线性约束在某点\emph{相切}、或起作用约束\emph{个数超过变量个数}时——这时务必检查 LICQ，别让一个失效的乘子方程把你引向错误结论。
}

\section{凸问题下 KKT 是充分条件}
\label{sec:kkt-suff}

到此为止 KKT 都只是\emph{必要}条件：最优点（在 CQ 下）必满足 KKT，但满足 KKT 的点未必最优——它可能是局部极小、鞍点，或诸多局部极大里的一个。这跟无约束情形里``一阶条件只给驻点''是同一个毛病。补上这道缺口的，仍然是凹凸性。

\thm{KKT 在凸问题下充分}{
考虑问题 $\max f(\vc x)$ s.t. $g_j(\vc x)\le0\ (j=1,\dots,m)$，并设
\begin{itemize}
    \item 目标 $f$ 是\emph{凹}函数；
    \item 每个约束函数 $g_j$ 是\emph{凸}函数（于是可行域 $\set{\vc x:g_j(\vc x)\le0,\forall j}$ 是凸集）。
\end{itemize}
若 $\vc x^*$ 连同某组乘子 $\vc\lambda^*\ge0$ 满足全部 KKT 条件 (i)--(iv)，则 $\vc x^*$ 是\textbf{全局}最优解。
}

\pf{
设 $\vc x^*,\vc\lambda^*$ 满足 KKT。构造 Lagrange 函数关于 $\vc x$ 的截面
\[
    \phi(\vc x)\;=\;f(\vc x)-\sum_{j}\lambda_j^*\,g_j(\vc x).
\]
因为 $f$ 凹、每个 $g_j$ 凸、且 $\lambda_j^*\ge0$，所以 $-\lambda_j^* g_j$ 凹，$\phi$ 作为凹函数之和仍\textbf{凹}。凹函数的一个基本性质是``切平面在其上方''（一阶条件，见第~\ref{ch:convex}~章）：对任意可行点 $\vc x$，
\[
    \phi(\vc x)\;\le\;\phi(\vc x^*)+\grad\phi(\vc x^*)\T(\vc x-\vc x^*).
\]
而 KKT 驻点条件 (i) 恰说 $\grad\phi(\vc x^*)=\grad f(\vc x^*)-\sum_j\lambda_j^*\grad g_j(\vc x^*)=\vc 0$，故上式右端的线性项消失：
\[
    \phi(\vc x)\;\le\;\phi(\vc x^*)\qquad\text{对一切 }\vc x .
\]
现把 $\phi$ 还原。任取一个\emph{可行}的 $\vc x$（即 $g_j(\vc x)\le0$）。由对偶可行 $\lambda_j^*\ge0$，有 $\sum_j\lambda_j^* g_j(\vc x)\le0$，因此
\[
    f(\vc x)\;=\;\phi(\vc x)+\sum_j\lambda_j^* g_j(\vc x)\;\le\;\phi(\vc x)\;\le\;\phi(\vc x^*).
\]
再看 $\vc x^*$ 这一头：由互补松弛 (iv)，$\sum_j\lambda_j^* g_j(\vc x^*)=0$，于是 $\phi(\vc x^*)=f(\vc x^*)$。两段拼起来，对一切可行 $\vc x$ 都有
\[
    f(\vc x)\;\le\;\phi(\vc x^*)\;=\;f(\vc x^*).
\]
即 $\vc x^*$ 是全局最大值点。
}

\state{凸性把 KKT 从``候选''升级为``答案''}{
这条定理是凸优化全部威力的来源，其分量与上一章``凹 $\To$ 局部最优即全局最优''完全平行：在凸问题（凹目标 + 凸可行域）里，KKT 条件\textbf{既必要又充分}。于是``解优化问题''彻底等价于``解 KKT 方程组''——找到任何一个满足四件套的 $(\vc x^*,\vc\lambda^*)$，就一劳永逸地拿到了全局最优，无需再担心局部解、无需比较多个候选。线性规划、二次规划、几乎所有现代凸优化求解器，骨子里都在求解 KKT 系统。此外凸问题里 Slater 条件（存在严格内点 $\bar{\vc x}$ 使所有 $g_j(\bar{\vc x})<0$）即可充当约束规范，比 LICQ 更易验证。
}

\rmkb[凸问题里 KKT 与对偶的接口]{
在凸问题中，KKT 条件还正好是原问题与其 Lagrange 对偶问题之间``强对偶''（对偶间隙为零）成立的刻画：$\vc x^*$ 原问题最优、$\vc\lambda^*$ 对偶问题最优，当且仅当二者满足 KKT。对偶视角（下一章「凸优化与对偶」）会把乘子 $\vc\lambda^*$ 进一步解读为对偶问题的解，并把影子价格的含义讲透。这里先记住：互补松弛正是连接原始与对偶两个世界的那座桥。
}

\section{实例：带预算与非负约束的消费选择}
\label{sec:kkt-example}

把四件套用一遍，最透亮的例子是消费者问题里的\textbf{角点解}——这也是 KKT 在微观经济学里最经典的落点。

\ex[角点解与互补松弛]{
消费者在两种商品上求效用极大，效用 $u(x_1,x_2)=\alpha\ln x_1+x_2$（$\alpha>0$，``拟线性''偏好），价格均为 $1$，收入为 $I>0$。问题是
\[
    \max_{x_1,x_2}\;\alpha\ln x_1+x_2
    \quad\text{s.t.}\quad
    x_1+x_2\le I,\;\; x_1\ge0,\;\; x_2\ge0 .
\]
用 KKT 求解，并说明收入 $I$ 很小时会出现什么。
}
\sol{
把约束统一写成 $g\le0$ 的形式：$g_1=x_1+x_2-I\le0$（预算），$g_2=-x_1\le0$，$g_3=-x_2\le0$（两条非负约束）。配乘子 $\lambda_1,\lambda_2,\lambda_3\ge0$，Lagrange 函数
\[
    \cL=\alpha\ln x_1+x_2-\lambda_1(x_1+x_2-I)-\lambda_2(-x_1)-\lambda_3(-x_2).
\]
驻点条件 (i)（对 $x_1,x_2$ 分别求偏导置零）：
\[
    \pd{\cL}{x_1}=\frac{\alpha}{x_1}-\lambda_1+\lambda_2=0,
    \qquad
    \pd{\cL}{x_2}=1-\lambda_1+\lambda_3=0 .
\]
效用对 $x_1$ 的边际 $\alpha/x_1$ 在 $x_1\to0$ 时趋于 $+\infty$（Inada 型），故最优必有 $x_1>0$；由互补松弛 $\lambda_2 x_1=0$ 得 $\lambda_2=0$。又因为多消费总有正效用（$x_2$ 的边际为 $1>0$），预算必花光，$g_1=0$ 即 $x_1+x_2=I$，对应 $\lambda_1>0$。剩下要判别的是 $x_2$ 这条非负约束的死活——这正是有没有角点解的关键。

\emph{情形一：内点解（$x_2>0$）。} 则 $\lambda_3=0$，第二式给 $\lambda_1=1$，第一式给 $\alpha/x_1=\lambda_1-\lambda_2=1$，即
\[
    x_1^*=\alpha,\qquad x_2^*=I-\alpha .
\]
要 $x_2^*>0$ 须 $I>\alpha$。这时拟线性偏好的标志性结论出现了：$x_1^*=\alpha$ 与收入\emph{无关}，收入的增减全部砸在 $x_2$ 上。

\emph{情形二：角点解（$x_2=0$）。} 当 $I\le\alpha$，上面的 $x_2^*=I-\alpha\le0$ 不可行，最优被逼到边界 $x_2=0$ 上。此时 $\lambda_3\ge0$ 不再为零，预算花光给 $x_1^*=I$，第一式（$\lambda_2=0$）给 $\lambda_1=\alpha/x_1=\alpha/I$，第二式给
\[
    \lambda_3=\lambda_1-1=\frac{\alpha}{I}-1 .
\]
当 $I<\alpha$ 时 $\lambda_3>0$，对偶可行满足，互补松弛 $\lambda_3 x_2=\lambda_3\cdot0=0$ 自动成立——这是一个合格的 KKT 点。它的含义是：收入太低，消费者把钱全花在商品 $1$ 上，商品 $2$ 一点不买，非负约束 $x_2\ge0$ 在此\textbf{紧绷}，其影子价格 $\lambda_3=\alpha/I-1>0$ 度量``如果允许 $x_2$ 取负值能多榨出多少效用''。

\emph{由于 $u$ 凹、约束全线性（自动凸、自动满足约束规范），由第~\ref{sec:kkt-suff}~节的充分性定理，上面找到的 KKT 点就是全局最优。}两种情形在 $I=\alpha$ 处平滑接上（$\lambda_3$ 从 $0$ 开始变正），互补松弛在此完成了从``$x_2$ 约束松''到``$x_2$ 约束紧''的切换。
}

图~\ref{fig:kkt-2} 把这两种情形并排画出，是理解互补松弛``切换''的最好图像。

\begin{figure}[h]
\centering
\begin{tikzpicture}[scale=0.92,>=Stealth]
  % ===== 左图：内点解 I>alpha (I=4, alpha=1.2) =====
  \begin{scope}
    \draw[->] (-0.2,0) -- (4.6,0) node[right] {$x_1$};
    \draw[->] (0,-0.2) -- (0,4.6) node[above] {$x_2$};
    % 预算线 x1+x2=I=4: 从(4,0)到(0,4)
    \fill[blue!8] (0,0) -- (4,0) -- (0,4) -- cycle;
    \draw[thick,blue!60!black] (4,0) -- (0,4);
    \node[blue!60!black,anchor=south west] at (2.25,1.95) {$x_1+x_2=I$};
    % 内点最优：x1=alpha=1.2, x2=I-alpha=2.8
    \fill[black] (1.2,2.8) circle (1.8pt);
    \node[anchor=south west] at (1.28,2.86) {$\vc x^*$};
    % 无差异曲线 u=alpha ln x1 + x2 过(1.2,2.8): u0=1.2*ln1.2+2.8=3.0188
    % 在最优处与预算线相切（斜率均为 -1）
    \draw[red!70!black,thick,domain=0.45:3.55,smooth,variable=\x]
        plot ({\x},{3.0188-1.2*ln(\x)});
    \node[red!70!black,anchor=west] at (2.4,3.45) {无差异曲线};
    \node[anchor=north] at (2.0,-0.35) {内点解：$I>\alpha$，$x_2^*>0$};
  \end{scope}
  % ===== 右图：角点解 I<alpha (I=2, alpha=2.6) =====
  \begin{scope}[xshift=6.6cm]
    \draw[->] (-0.2,0) -- (4.6,0) node[right] {$x_1$};
    \draw[->] (0,-0.2) -- (0,4.6) node[above] {$x_2$};
    % 预算线 x1+x2=I=2: 从(2,0)到(0,2)
    \fill[blue!8] (0,0) -- (2,0) -- (0,2) -- cycle;
    \draw[thick,blue!60!black] (2,0) -- (0,2);
    \node[blue!60!black,anchor=north east] at (1.35,0.62) {$x_1+x_2=I$};
    % 角点最优：x1=I=2, x2=0
    \fill[black] (2,0) circle (2.0pt);
    \node[anchor=south west] at (2.08,0.12) {$\vc x^*$};
    % 无差异曲线过(2,0)：u0=2.6 ln(2)=1.8022; x2=1.8022-2.6 ln x1
    % 在 x1=2 处斜率 -alpha/x1=-1.3，比预算线 -1 更陡 => 角点；曲线在 x1=2 处恰交 x2=0
    \draw[red!70!black,thick,domain=0.62:2.0,smooth,variable=\x]
        plot ({\x},{1.8022-2.6*ln(\x)});
    \node[red!70!black,anchor=west] at (2.3,2.75) {无差异曲线};
    \draw[red!70!black,->] (2.25,2.6) -- (1.05,1.9); % 引线指向曲线上方段，不穿预算线/标签
    \node[anchor=north] at (2.0,-0.35) {角点解：$I<\alpha$，$x_2^*=0$};
  \end{scope}
\end{tikzpicture}
\caption{互补松弛的切换。左：收入较高（$I>\alpha$），无差异曲线与预算线在内部相切，两种商品都买，非负约束 $x_2\ge0$ 松弛、$\lambda_3=0$。右：收入较低（$I<\alpha$），在边界 $x_2=0$ 处无差异曲线比预算线更陡（边际替代率始终大于价格比），最优被顶到角点，$x_2\ge0$ 紧绷、影子价格 $\lambda_3=\alpha/I-1>0$。从左到右，正是互补松弛把约束 $x_2\ge0$ 由``松''切到``紧''。}
\label{fig:kkt-2}
\end{figure}

\section{它通向何处}
\label{sec:kkt-beyond}

KKT 是带约束优化的通用语言，本章这把工具的去处密布在三大核心课乃至应用领域里：

\begin{itemize}
    \item \textbf{微观经济学：角点解与非负约束}。消费者理论里``某商品买零''、生产理论里``某要素不投入''、以及一切``非负''变量，都是非负约束起作用的角点解，全靠互补松弛刻画。上例的拟线性需求只是冰山一角。
    \item \textbf{机制设计与契约理论：IC 与 IR 约束}。激励相容（IC）与参与（IR）约束本质都是不等式约束，最优契约的特征（哪条 IC 紧、哪条 IR 紧）正是一组互补松弛的结论——``哪类代理人的约束绑定''往往就是整个机制的设计要点。
    \item \textbf{凸优化与对偶}（下一章）。凸问题下 KKT 充分性 + Slater 约束规范，是整个凸优化与 Lagrange 对偶理论的支点；线性规划、二次规划都是其特例。
    \item \textbf{受约束的估计：受限 MLE / GMM}。当参数被理论先验地限制（方差非负、概率落在单纯形上、份额加总为一、单调性或不等式形状约束），极大似然或 GMM 就变成带不等式约束的优化，其解由 KKT 刻画；约束\emph{是否起作用}还直接影响检验统计量的渐近分布（边界上的参数会让标准正态近似失效，这是``边界参数推断''问题的根源，渐近统计部分会再遇到）。
\end{itemize}

\state{一句话收束}{
Lagrange 法回答``约束恰好取等时怎么办'',KKT 把它推广到``约束可松可紧''的真实世界，代价只是多了一组对偶可行（$\lambda\ge0$）与互补松弛（$\lambda_j g_j=0$）的条件——而后者恰恰自动解决了``哪些约束起作用''这个最棘手的问题。再叠上凹凸性，KKT 就从``最优的必要特征''升格为``找到全局最优的充要判据''，成为凸优化与一切受约束决策、受约束估计的共同地基。
}
